\documentclass[11pt,a4paper]{article}

% ===== PACKAGES =====
\usepackage[utf8]{inputenc}
\usepackage[T1]{fontenc}
\usepackage[margin=1in]{geometry}
\usepackage{hyperref}
\usepackage{listings}
\usepackage{xcolor}
\usepackage{booktabs}
\usepackage{array}
\usepackage{longtable}
\usepackage{fancyhdr}
\usepackage{titlesec}
\usepackage{tcolorbox}
\usepackage{fontawesome5}
\usepackage{amssymb}
\usepackage{amsmath}
\usepackage{enumitem}

% ===== COLORS =====
\definecolor{codegreen}{rgb}{0.25,0.49,0.48}
\definecolor{codegray}{rgb}{0.5,0.5,0.5}
\definecolor{codepurple}{rgb}{0.58,0,0.82}
\definecolor{codeblue}{rgb}{0.13,0.29,0.53}
\definecolor{codeorange}{rgb}{0.8,0.4,0.0}
\definecolor{backcolour}{rgb}{0.95,0.95,0.95}
\definecolor{warningbg}{rgb}{1.0,0.95,0.85}
\definecolor{warningborder}{rgb}{0.9,0.6,0.2}
\definecolor{tipbg}{rgb}{0.9,0.97,0.9}
\definecolor{tipborder}{rgb}{0.3,0.7,0.3}

% ===== CODE LISTING STYLE =====
\lstdefinelanguage{IEC61131}{
    keywords={VAR, END_VAR, VAR_INPUT, VAR_OUTPUT, VAR_IN_OUT, VAR_GLOBAL, 
              TYPE, END_TYPE, STRUCT, END_STRUCT, FUNCTION, END_FUNCTION,
              FUNCTION_BLOCK, END_FUNCTION_BLOCK, PROGRAM, END_PROGRAM,
              IF, THEN, ELSE, ELSIF, END_IF, CASE, OF, END_CASE,
              FOR, TO, BY, DO, END_FOR, WHILE, END_WHILE, REPEAT, UNTIL, END_REPEAT,
              TRUE, FALSE, AND, OR, NOT, XOR, MOD,
              ARRAY, OF, POINTER, REFERENCE, TO, AT},
    keywordstyle=\color{codeblue}\bfseries,
    keywords=[2]{BOOL, BYTE, WORD, DWORD, LWORD, SINT, USINT, INT, UINT, 
                 DINT, UDINT, LINT, ULINT, REAL, LREAL, STRING, WSTRING,
                 TIME, LTIME, DATE, TIME_OF_DAY, TOD, DATE_AND_TIME, DT,
                 TcEventSeverity, T_FILETIME},
    keywordstyle=[2]\color{codepurple}\bfseries,
    keywords=[3]{ADR, SIZEOF, REF, FILETIME_TO_DT, GETSYSTEMTIME},
    keywordstyle=[3]\color{codeorange}\bfseries,
    sensitive=false,
    morecomment=[l]{//},
    morecomment=[s]{(*}{*)},
    commentstyle=\color{codegreen}\itshape,
    stringstyle=\color{codeorange},
    morestring=[b]',
    morestring=[b]",
}

\lstdefinestyle{iecstyle}{
    language=IEC61131,
    backgroundcolor=\color{backcolour},
    basicstyle=\ttfamily\small,
    breakatwhitespace=false,
    breaklines=true,
    captionpos=b,
    keepspaces=true,
    numbers=left,
    numbersep=8pt,
    numberstyle=\tiny\color{codegray},
    showspaces=false,
    showstringspaces=false,
    showtabs=false,
    tabsize=4,
    frame=single,
    framerule=0.5pt,
    rulecolor=\color{codegray},
    xleftmargin=15pt,
    framexleftmargin=15pt,
}

\lstset{style=iecstyle}

% ===== TCOLORBOX STYLES =====
\tcbuselibrary{skins,breakable}

\newtcolorbox{warningbox}{
    colback=warningbg,
    colframe=warningborder,
    boxrule=1pt,
    left=10pt,
    right=10pt,
    top=8pt,
    bottom=8pt,
    breakable,
    before upper={\faExclamationTriangle\ \textbf{Warning:} }
}

\newtcolorbox{tipbox}{
    colback=tipbg,
    colframe=tipborder,
    boxrule=1pt,
    left=10pt,
    right=10pt,
    top=8pt,
    bottom=8pt,
    breakable,
    before upper={\faLightbulb\ \textbf{Tip:} }
}

% ===== HEADER/FOOTER =====
\pagestyle{fancy}
\fancyhf{}
\fancyhead[L]{\leftmark}
\fancyhead[R]{TwinCAT 3 Lecture Notes}
\fancyfoot[C]{\thepage}
\renewcommand{\headrulewidth}{0.4pt}
\renewcommand{\footrulewidth}{0.4pt}

% ===== HYPERREF SETUP =====
\hypersetup{
    colorlinks=true,
    linkcolor=codeblue,
    filecolor=magenta,
    urlcolor=cyan,
    pdftitle={TwinCAT 3 - Structures and Functions},
    pdfauthor={},
    bookmarks=true,
}

% ===== TITLE FORMATTING =====
\titleformat{\section}
    {\normalfont\Large\bfseries}{\thesection}{1em}{}[{\titlerule[0.8pt]}]
\titleformat{\subsection}
    {\normalfont\large\bfseries}{\thesubsection}{1em}{}
\titleformat{\subsubsection}
    {\normalfont\normalsize\bfseries}{\thesubsubsection}{1em}{}

% ===== DOCUMENT INFO =====
\title{
    \vspace{-1cm}
    \textbf{Lecture Notes: TwinCAT 3}\\
    \Large Structures and Functions
}
\author{}
\date{\today}

% ===== BEGIN DOCUMENT =====
\begin{document}

\maketitle

\tableofcontents
\newpage

% ============================================================
\section{Introduction to Structures (DUTs)}
% ============================================================

While \textbf{arrays} are used to hold multiple data items of the same kind, \textbf{structures} allow you to hold several data items of \textbf{different kinds}. Structures are user-defined data types (UDT) that represent a record or a logical group.

\begin{itemize}[leftmargin=*]
    \item \textbf{Advantages:} They are more user-friendly than maintaining multiple separate arrays for related data. If a data definition changes, you only need to update it in one place (the structure definition) rather than everywhere those variables are used.
    
    \item \textbf{Syntax:} To access components within a structure, use \textbf{dot notation} (e.g., \texttt{StructureName.ComponentName}).
    
    \item \textbf{Creation:} Right-click on the \textbf{DUTs} folder in the project tree, select \textbf{Add > DUT}, choose \textbf{Structure}, and give it a name (often prefixed with \texttt{ST\_} by convention).
    
    \item \textbf{Scope:} In TwinCAT 3, structures have a \textbf{global scope}, meaning they are visible to all code within the project.
\end{itemize}

\subsection*{Example: Structure Definition}

\begin{lstlisting}
// DUT Definition: ST_Event
// Created as a separate DUT file in the project tree under the DUTs folder
TYPE ST_Event :
STRUCT
    eEventType     : E_EventType;        // User-defined ENUM for event classification
    eEventSeverity : TcEventSeverity;    // Beckhoff event severity type
    nIdentity      : UDINT;              // Unique identifier for the event
    sText          : STRING(255);        // Descriptive text for the event
    dtTimestamp    : DATE_AND_TIME;      // When the event occurred
END_STRUCT
END_TYPE
\end{lstlisting}

\begin{tipbox}
Use the \texttt{ST\_} prefix convention for structure names to easily identify them throughout your codebase and distinguish them from other data types.
\end{tipbox}

% ============================================================
\section{Functions (POUs)}
% ============================================================

Functions are essential for \textbf{modularization and code reuse}; you define the logic once and can use it many times throughout your software.

\begin{itemize}[leftmargin=*]
    \item \textbf{Components:} A function consists of a name, an optional \textbf{return type} (e.g., REAL, INT), \textbf{input parameters}, and \textbf{local variables} that act as a ``scratch pad''.
    
    \item \textbf{Returning Values:} To return a value from a function, you assign the result directly to the \textbf{function's name} within the code.
    
    \item \textbf{Advanced Features:}
    \begin{itemize}
        \item \textbf{Multiple Outputs:} Functions can be designed to return more than one value by defining additional output parameters.
        \item \textbf{VAR\_IN\_OUT:} This keyword allows variables to act as both inputs and outputs, meaning the function can both read from and write to the same variable.
    \end{itemize}
\end{itemize}

\subsection*{Example: Basic Function}

\begin{lstlisting}
// Function: Celsius_To_Fahrenheit
// Converts a temperature value from Celsius to Fahrenheit
FUNCTION Celsius_To_Fahrenheit : REAL
VAR_INPUT
    fCelsius : REAL;    // Input temperature in Celsius
END_VAR
VAR
    // Local variables (scratch pad) - none needed for this simple function
END_VAR

// Calculate and return the Fahrenheit value
// Assign the result to the function name to return it
Celsius_To_Fahrenheit := fCelsius * 1.8 + 32.0;
\end{lstlisting}

\subsubsection*{Usage Example}

\begin{lstlisting}
// Calling the function from a program or function block
VAR
    fTempCelsius    : REAL := 25.0;
    fTempFahrenheit : REAL;
END_VAR

// Call the function and store the result
fTempFahrenheit := Celsius_To_Fahrenheit(fCelsius := fTempCelsius);
// Result: fTempFahrenheit = 77.0
\end{lstlisting}

% ============================================================
\section{Passing Parameters: Value vs. Reference}
% ============================================================

When sending data into a function, it can be passed in two ways:

\begin{itemize}[leftmargin=*]
    \item \textbf{Pass by Value:} The function receives a \textbf{copy} of the variable. Changes made inside the function do not affect the original data. This is best for small, simple data types like INT, REAL, or BOOL.
    
    \item \textbf{Pass by Reference:} The function receives the \textbf{memory location} of the original variable. Any changes made inside the function persist after the function completes. This is more efficient for large data sets like \textbf{structures} because it avoids the ``memory penalty'' of copying large amounts of data.
    
    \item \textbf{Implementation:} In TwinCAT, passing by reference is achieved using either the \texttt{VAR\_IN\_OUT} keyword or by defining a parameter as a \texttt{REFERENCE TO} a specific type.
\end{itemize}

\begin{table}[h]
\centering
\begin{tabular}{@{}lll@{}}
\toprule
\textbf{Aspect} & \textbf{Pass by Value} & \textbf{Pass by Reference} \\
\midrule
Data Passed & Copy of the variable & Memory address \\
Original Affected & No & Yes \\
Best For & Small types (INT, REAL, BOOL) & Large types (structures, arrays) \\
Memory Overhead & Higher (copy created) & Lower (no copy) \\
Keywords & \texttt{VAR\_INPUT} & \texttt{VAR\_IN\_OUT}, \texttt{REFERENCE TO} \\
\bottomrule
\end{tabular}
\caption{Comparison of Parameter Passing Methods}
\end{table}

\subsection*{Example: Passing by Reference}

\begin{lstlisting}
// Function: UpdateEventTimestamp
// Updates the timestamp of an event structure passed by reference
FUNCTION UpdateEventTimestamp : BOOL
VAR_INPUT
    stEvent : REFERENCE TO ST_Event;    // Reference to the event structure
END_VAR
VAR
    fbGetSystemTime : GETSYSTEMTIME;    // Function block instance for system time
    stFileTime      : T_FILETIME;       // Temporary storage for file time
END_VAR

// Get the current system time
fbGetSystemTime(timeLoDW => stFileTime.dwLowDateTime,
                timeHiDW => stFileTime.dwHighDateTime);

// Directly modify the referenced structure's timestamp member
// This change persists after the function returns
stEvent.dtTimestamp := FILETIME_TO_DT(stFileTime);

// Return TRUE to indicate successful update
UpdateEventTimestamp := TRUE;
\end{lstlisting}

\begin{warningbox}
When passing by reference, any modifications made to the parameter inside the function will directly affect the original variable. Ensure this is the intended behavior before using \texttt{REFERENCE TO} or \texttt{VAR\_IN\_OUT}.
\end{warningbox}

% ============================================================
\section{Practical Implementation: System Time \& Libraries}
% ============================================================

The sources illustrate these concepts through a practical example: a function that updates an event's timestamp using the PLC's system time.

\begin{itemize}[leftmargin=*]
    \item \textbf{External Libraries:} To access advanced functionality like system time, you may need to add references to libraries like \textbf{Tc2\_Utilities}.
    
    \item \textbf{Format Conversion:} The system time is often provided as a 64-bit integer (FILETIME) which must be converted into a human-readable \texttt{DATE\_AND\_TIME} format using built-in functions like \texttt{FILETIME\_TO\_DT}.
    
    \item \textbf{Function Blocks (Preview):} To get the system time, the example uses a \textbf{Function Block} (\texttt{GETSYSTEMTIME}), which must be \textbf{instantiated} in the local variables before use.
\end{itemize}

\subsection*{Example: System Time Conversion}

\begin{lstlisting}
// Declaration section
VAR
    fbGetSystemTime : GETSYSTEMTIME;    // Instance of the system time function block
    stFileTime      : T_FILETIME;       // Structure to hold the raw file time
    dtCurrentTime   : DATE_AND_TIME;    // Human-readable date and time result
END_VAR

// Implementation section
// Step 1: Call the GETSYSTEMTIME function block to retrieve system time
// The function block outputs the time as two 32-bit values (low and high DWORD)
fbGetSystemTime(
    timeLoDW => stFileTime.dwLowDateTime,   // Lower 32 bits of file time
    timeHiDW => stFileTime.dwHighDateTime   // Upper 32 bits of file time
);

// Step 2: Convert the T_FILETIME structure to a DATE_AND_TIME value
// FILETIME_TO_DT is a built-in conversion function from Tc2_Utilities
dtCurrentTime := FILETIME_TO_DT(stFileTime);

// dtCurrentTime now contains the current date and time in a readable format
// Example value: DT#2024-06-15-14:30:45
\end{lstlisting}

\subsection{Adding the Tc2\_Utilities Library}

To use the \texttt{GETSYSTEMTIME} function block and \texttt{FILETIME\_TO\_DT} function, you must first add the library reference:

\begin{enumerate}
    \item In the Solution Explorer, expand your PLC project
    \item Right-click on \textbf{References}
    \item Select \textbf{Add library...}
    \item Search for \textbf{Tc2\_Utilities} and add it
\end{enumerate}

\begin{tipbox}
The \texttt{T\_FILETIME} structure represents time as the number of 100-nanosecond intervals since January 1, 1601. The \texttt{FILETIME\_TO\_DT} function handles this conversion automatically, making it easy to work with human-readable timestamps.
\end{tipbox}

\subsection{Complete Working Example}

Below is a complete example combining all the concepts from this lecture:

\begin{lstlisting}
// Function: LogEvent
// Creates and logs an event with the current system timestamp
FUNCTION LogEvent : BOOL
VAR_INPUT
    stEvent      : REFERENCE TO ST_Event;   // Event to be logged (by reference)
    eType        : E_EventType;             // Type of event
    eSeverity    : TcEventSeverity;         // Severity level
    nId          : UDINT;                   // Event identifier
    sDescription : STRING(255);             // Event description
END_VAR
VAR
    fbGetSystemTime : GETSYSTEMTIME;        // System time function block instance
    stFileTime      : T_FILETIME;           // Temporary file time storage
END_VAR

// Populate the event structure members
stEvent.eEventType     := eType;
stEvent.eEventSeverity := eSeverity;
stEvent.nIdentity      := nId;
stEvent.sText          := sDescription;

// Get and set the current timestamp
fbGetSystemTime(timeLoDW => stFileTime.dwLowDateTime,
                timeHiDW => stFileTime.dwHighDateTime);
stEvent.dtTimestamp := FILETIME_TO_DT(stFileTime);

// Return success
LogEvent := TRUE;
\end{lstlisting}

% ============================================================
\section*{Quick Reference Card}
\addcontentsline{toc}{section}{Quick Reference Card}
% ============================================================

\subsection*{Structure Declaration Syntax}

\begin{lstlisting}
TYPE <StructureName> :
STRUCT
    <member1> : <DataType1>;
    <member2> : <DataType2>;
    // ... additional members
END_STRUCT
END_TYPE
\end{lstlisting}

\subsection*{Function Declaration Syntax}

\begin{lstlisting}
FUNCTION <FunctionName> : <ReturnType>
VAR_INPUT
    <inputParam> : <DataType>;              // Pass by value
    <refParam>   : REFERENCE TO <DataType>; // Pass by reference
END_VAR
VAR_IN_OUT
    <inOutParam> : <DataType>;              // Pass by reference (read/write)
END_VAR
VAR
    <localVar> : <DataType>;                // Local scratch variables
END_VAR

// Function body
<FunctionName> := <returnValue>;            // Assign return value
\end{lstlisting}

\subsection*{Common Naming Conventions}

\begin{table}[h]
\centering
\begin{tabular}{@{}lll@{}}
\toprule
\textbf{Prefix} & \textbf{Type} & \textbf{Example} \\
\midrule
\texttt{ST\_} & Structure & \texttt{ST\_Event}, \texttt{ST\_MotorData} \\
\texttt{E\_} & Enumeration & \texttt{E\_EventType}, \texttt{E\_State} \\
\texttt{fb} & Function Block Instance & \texttt{fbGetSystemTime} \\
\texttt{st} & Structure Instance & \texttt{stEvent}, \texttt{stFileTime} \\
\texttt{dt} & DATE\_AND\_TIME & \texttt{dtTimestamp}, \texttt{dtCurrentTime} \\
\bottomrule
\end{tabular}
\caption{Common TwinCAT Naming Conventions}
\end{table}

\vfill

\begin{center}
\textit{Last Updated: \today}
\end{center}

\end{document}